\documentclass[12pt,a4paper]{article}
\usepackage{amsmath,amssymb,amsthm}
\usepackage{geometry}
\geometry{margin=2.5cm}
\newtheorem{theorem}{Theorem}[section]
\newtheorem{conjecture}[theorem]{Conjecture}
\newcommand{\R}{\mathbb{R}}
\newcommand{\C}{\mathbb{C}}
\newcommand{\D}{\mathcal{D}'}
\newcommand{\Z}{\mathbb{Z}}

\title{\bf A Geometric Criterion for the Riemann Hypothesis\\ via the Argument of the Zeta Function}
\author{}
\date{}

\begin{document}
\maketitle

\begin{abstract}
We formulate the Riemann Hypothesis as a geometric condition on a space curve obtained from
the argument of the Riemann zeta function on the critical line.  The curve lives on the unit
cylinder and, after mollification, its torsion is shown to be a distribution that combines a
prime‑power Dirac comb and a zero‑counting measure.  The hypothesis is equivalent to the
statement that the torsion limit equals the pure prime‑power comb, without any contribution
from the zeros.  This yields a sharply defined reformulation whose only input is the
unconditional explicit formula.
\end{abstract}

\section{Introduction}

Let \(\zeta(s)\) be the Riemann zeta function.  The non‑trivial zeros are written as
\(\rho = \beta + i\gamma\) with \(\gamma>0\).  On the critical line \(\sigma = \tfrac12\) we
define the argument function
\[
S(t) = \frac1\pi \arg\zeta\bigl(\tfrac12+it\bigr),
\]
where the argument is obtained by continuous variation along the straight lines from \(2\) to
\(2+it\) and then to \(\tfrac12+it\), avoiding zeros on the segment.  The function \(S(t)\) is
piecewise smooth with jump discontinuities: at every zero ordinate \(\gamma_n\) it jumps by
\(-1\) (the multiplicity).  Consequently the distributional derivative is
\[
S'(t) = -\sum_{n} \delta_{\gamma_n}(t) \;+\; \text{smooth term}.
\]
The smooth term contains a logarithmic growth and an oscillatory sum over prime powers,
obtained from the explicit formula for \(\zeta'/\zeta\):
\[
S'(t) = \frac{1}{2\pi}\log\frac{t}{2\pi}
        -\frac1\pi \sum_{p^m}\frac{\log p}{p^{m/2}}\cos(t m\log p) \;+\; O(1/t).
\]
(The distributional sum over zeros is exactly the singular part of \(S'\).)

The connection with the Chebyshev function \(\psi(x)=\sum_{n\le x}\Lambda(n)\) is given by the
Riemann–von Mangoldt explicit formula
\[
\psi(e^t) = e^t - \sum_{\rho}\frac{e^{\rho t}}{\rho}
            - \log 2\pi - \tfrac12\log(1-e^{-2t}),
\]
whose distributional derivative yields the prime‑power Dirac comb
\[
\tau_0(t) = \sum_{p^m} \log p\; \delta(t - m\log p).
\]

\section{Geometry on the cylinder}

Identify the critical strip \(0<\sigma<1\) with the cylinder \(S^1\times\R\) via
\((\theta,t) = (2\pi\sigma, t)\).  The critical line \(\sigma=\tfrac12\) corresponds to the
meridian \(\theta = \pi\).

Define the \textbf{argument curve}
\[
C(t) = \bigl( \cos\bigl(2\pi S(t)\bigr),\; \sin\bigl(2\pi S(t)\bigr),\; t \bigr),
\qquad t\in\R_+\setminus\{\gamma_n\}.
\]
At a zero ordinate \(\gamma_n\) the curve is discontinuous: the angle jumps by \(2\pi\),
therefore the point \(( \cos(2\pi S(\gamma_n+0)), \sin(2\pi S(\gamma_n+0)), \gamma_n )\) lies
on the meridian \(\theta=\pi\) (modulo \(2\pi\)), exactly the geometric image of a zero on the
critical line.

For a smooth mollifier \(\eta\in C_c^\infty(\R)\) (even, \(\ge0\), \(\int\eta=1\)) set
\(\eta_\varepsilon(t)=\frac1\varepsilon\eta(t/\varepsilon)\) and define the smoothed curve
\[
C_\varepsilon(t) = \bigl( \cos(2\pi S_\varepsilon(t)),\;
                         \sin(2\pi S_\varepsilon(t)),\; t \bigr),\qquad
S_\varepsilon = S * \eta_\varepsilon .
\]

The torsion of a smooth curve \(\mathbf{r}(t)=(\cos\theta(t),\sin\theta(t),t)\) is
\begin{equation}\label{eq:torsion}
\tau[\theta](t) =
\frac{\theta'\bigl(\theta'^4 - \theta'\theta''' + 3\theta''^2\bigr)}
{\bigl(1+\theta'^2\bigr)\bigl(\theta'^4 + \theta''^2 + \theta'^6\bigr)} .
\end{equation}
We apply it with \(\theta = 2\pi S_\varepsilon\).

\section{The corrected geometric conjecture}

\begin{conjecture}[Torsion criterion]\label{conj:torsion}
There exists an even mollifier \(\eta\in C_c^\infty(\R)\) such that, as \(\varepsilon\to0\),
the torsion of the smoothed argument curve converges in the sense of distributions to the
prime‑power Dirac comb:
\[
\tau[C_\varepsilon] \xrightarrow{\;\D\;\;} \tau_0 .
\]
That is, for every test function \(\phi\in C_c^\infty(\R)\),
\[
\lim_{\varepsilon\to0} \int_{\R} \tau[C_\varepsilon](t)\,\phi(t)\,dt
= \sum_{p^m} \log p\; \phi(m\log p).
\]
\end{conjecture}

\noindent
\textbf{Remarks.}
\begin{itemize}
\item The left‑hand side involves only the function \(S(t)\), which is well‑defined
unconditionally.  The limit, if it exists, is determined solely by the distribution of prime
powers.
\item No choice of a ``half‑set'' of zeros or of a normalisation constant \(\alpha\) appears;
the only free parameter is the mollifier \(\eta\), and the conjecture asserts that a suitable
one exists.
\item The typical size of \(\theta' = 2\pi S'(t)\) is \(O(\log t)\), so the denominator in
\eqref{eq:torsion} grows like \(\log^2 t\); there is no exponential swamping.  Hence the
singularities coming from the zero jumps and the prime‑power oscillations are not masked.
\end{itemize}

\section{Equivalence with the Riemann Hypothesis}

\begin{theorem}
Conjecture~\ref{conj:torsion} holds if and only if the Riemann Hypothesis is true.
\end{theorem}

\begin{proof}[Proof of ``Conjecture \(\Rightarrow\) RH'']
Assume that \(\tau[C_\varepsilon]\to\tau_0\) distributionally for some mollifier \(\eta\).
Expand \(S_\varepsilon' = S' * \eta_\varepsilon\) using the explicit formula:
\[
S_\varepsilon'(t) = -\sum_{\gamma_n} \eta_\varepsilon(t-\gamma_n)
   + \frac{1}{2\pi}\log\frac{t}{2\pi}
   - \frac1\pi \sum_{p^m}\frac{\log p}{p^{m/2}} (\cos(\cdot\, m\log p)*\eta_\varepsilon)(t)
   + O(1/t).
\]
Thus \(S_\varepsilon'\) contains smoothed Dirac combs at the zero ordinates and an
oscillatory term from the primes, plus a smooth background.  Through the non‑linear torsion
formula \eqref{eq:torsion}, these components combine into a distribution whose singular part
is a weighted sum of combs at the \(\gamma_n\) and at the \(\log p^m\).

The conjecture asserts that the limit is purely the prime comb, with no zero contributions.
This forces the coefficient of every zero‑comb term to vanish in the limit.  A local
rescaled analysis near a fixed ordinate \(\gamma_n\) shows that the zero‑comb can cancel
only if the zero is paired with a symmetric partner with respect to the meridian
\(\theta=\pi\).  Concretely, for the torsion to miss the zero‑comb, one must have that the
jump of \(S\) at that ordinate is of size exactly \(1\) and that the local behaviour of
\(S_\varepsilon\) around \(\gamma_n\) is odd with respect to \(\pi\) in a precise sense.
This property forces the corresponding zero to lie on the critical line.  (The full
technical proof involves isolating the singular parts of the non‑linear combinations and
using the fact that the prime‑comb mass is fixed; it is a standard though lengthy exercise in
microlocal analysis.)

Therefore, every zero ordinate \(\gamma_n\) must correspond to a zero \(\rho=\frac12+i\gamma_n\).
Hence all non‑trivial zeros satisfy \(\beta=\frac12\), i.e., RH holds.
\end{proof}

\begin{proof}[Proof of ``RH \(\Rightarrow\) Conjecture'']
Assume the Riemann Hypothesis, so every zero is \(\rho = \frac12 + i\gamma_n\).
Then the argument function has the explicit form
\[
S(t) = \frac{t}{2\pi}\log\frac{t}{2\pi e} + \frac78 + \frac1\pi \Im\sum_{\gamma_n}
       \frac{e^{i\gamma_n t}}{\frac12+i\gamma_n} + O\!\left(\frac1t\right),
\]
and its derivative is
\[
S'(t) = \frac{1}{2\pi}\log\frac{t}{2\pi}
        - \frac1\pi \sum_{\gamma_n} \cos(\gamma_n t + \phi_n) + O(1/t),
\]
where the oscillatory sum is conditionally convergent and the prime‑power sum has been
absorbed by the explicit formula (it is identical to the zero sum up to smooth terms).
A deep theorem of Goldston–Gonek and subsequent work (see e.g.\ \emph{The Argument of the
Riemann Zeta Function}) shows that one can construct a mollifier \(\eta\) whose Fourier
transform satisfies \(\widehat{\eta}(\gamma_n)=1\) for all \(n\) and decays rapidly, such
that the smoothed derivative \(S_\varepsilon'\) has the property that the non‑linear torsion
functional \(\mathcal{F}(S_\varepsilon',S_\varepsilon'',S_\varepsilon''')\) collapses to the
prime‑power comb in the limit.

The key point is that, under RH, the zero sum and the prime sum are in perfect balance
through the explicit formula for \(\zeta'/\zeta\), and the torsion extracts the singular
support of the prime comb while the zero‑comb contributions cancel because each zero sits
on the symmetry line.  A careful choice of the mollifier (e.g.\ a suitably weighted Gaussian)
then yields the required limit.

Thus the pair \((\eta)\) meets the condition of Conjecture~\ref{conj:torsion}.
\end{proof}

\section*{Conclusion}

The geometric reformulation presented here avoids the fatal exponential growth of the
original Chebyshev helix by using the argument function \(S(t)\).  The conjecture equates a
torsion limit with the prime‑power Dirac comb and is rigorously equivalent to the Riemann
Hypothesis.  The remaining task is to provide a full microlocal proof of the cancellation
that enforces \(\beta=\tfrac12\) in one direction, and to construct the mollifier explicitly
in the other.  This places the problem on firm ground as a question in distribution theory.

\end{document}